% !TeX program = lualatex
% !TeX encoding = utf8
% !TeX spellcheck = uk_UA
% !TEX indexprogram = upmendex
% !TeX root =./LabWork_Part1.tex

%=========================================================
\chapter*{Передмова}\label{\currfilebase}
\addcontentsline{toc}{chapter}{Передмова}
%=========================================================


Даний навчальний посібник є першою частиною лабораторного практикуму
з курсу «Електрика та магнетизм», який читається студентам спеціальності
E6~«Прикладна фізика та наноматеріали». Посібник охоплює
розділи пов'язані з електростатичним і
магнітостатичним полем та постійним струмом, і призначений для
закріплення теоретичних знань, отриманих на лекційних та практичних
заняттях, шляхом безпосередньої експериментальної перевірки
фундаментальних законів електромагнетизму.

В посібнику наведено п'ять лабораторних робіт. Перші дві роботи
присвячені дослідженню електростатичного поля: перевірці закону Кулона
за допомогою крутильних терезів та вивченню розподілу напруженості й
потенціалу заряджених тіл методом потенціального зонда. Третя робота
знайомить студентів із характеристиками джерел електричної енергії~---
електрорушійною силою, внутрішнім опором та режимами роботи кола.
Четверта та п'ята роботи присвячені магнітостатиці: визначенню
магнітного моменту контуру зі струмом у полі котушок Гельмгольца та
експериментальному вимірюванню компонент магнітного поля Землі.

Кожна робота побудована за єдиною структурою й містить: формулювання
мети та ключових понять, стислий виклад теоретичних відомостей,
необхідних для свідомого виконання експерименту, опис ідеї досліду та
обладнання, покроковий хід роботи, а також перелік контрольних питань
і розрахункових завдань для самостійного опрацювання. Такий підхід
дозволяє студенту не лише механічно виконати вимірювання, а й
зрозуміти фізичну суть явищ, що досліджуються, оцінити межі
застосовності теоретичних моделей та навчитися коректно аналізувати
похибки експерименту.

Особливу увагу в посібнику приділено історичному та методологічному
контексту: наводяться відомості про розвиток уявлень про електричний
заряд і магнетизм, огляд класичних і сучасних експериментальних
перевірок закону Кулона, а також обговорення сучасних гіпотез щодо
походження геомагнітного поля. Це покликано сформувати в студентів не
лише практичні навички роботи з вимірювальною апаратурою, а й цілісне
розуміння місця експериментальної фізики в розвитку природничих наук.

\bigskip
\noindent
\hfill\textit{Автори}